%%%%%%%%%%  scan idx 552  |  folio 539  %%%%%%%%%%
\SGAPrintPage{539}%

\(\Grb(X)\) thus becomes a module over the ring \(\Gro(X)\) (and even a graded module over the graded ring \(\Gro(X)\), provided the sign of the grading of \(\Grb(X)\) is changed by putting \(\Gr^i_\bullet(X)=\Gr_{-i}(X)\) \((i\in\mathbb Z)\)).
Now consider a proper morphism
\[
 f:X'\longrightarrow X
\]
of noetherian schemes. Then \(f\) induces a homomorphism
\[
 f^*:\Gro(X)\longrightarrow\Gro(X')
\]
of graded rings (V), and, by 1.1.8, a homomorphism
\[
 f_*:\Grb(X')\longrightarrow\Grb(X)
\]
of graded abelian groups. Passing to the associated graded groups, the projection formula (IV 2.11) and (3.1.1) then give the graded projection formula
\sgaStableTarget{sga6.e10.eq.3-1-2}\begin{equation}\tag{3.1.2}
 x\cdot f_*(y)=f_*\bigl(f^*(x)\cdot y\bigr)
 \qquad (x\in\Gro(X),\ y\in\Grb(X')).
\end{equation}

\sgaStableTarget{sga6.e10.sec.3.2}\subsection*{3.2.}
Let \(f:X'\to X\) be a projective complete-intersection morphism (Exp. VIII) of virtual dimension \(r\) between noetherian schemes, with \(X\) carrying an ample invertible \(\underline{O}_X\)-module. In Exp. VIII it was proved that \(f\) induces a homomorphism
\sgaStableTarget{sga6.e10.eq.3-2-1}\begin{equation}\tag{3.2.1}
 f_*:\Filt^k(X')_{\mathbb Q}\longrightarrow\Filt^{k-r}(X)_{\mathbb Q}\qquad(k\in\mathbb Z)
\end{equation}
which moreover satisfies the projection formula (Exp. IV):
\sgaStableTarget{sga6.e10.eq.3-2-2}\begin{equation}\tag{3.2.2}
 f_*(f^*(x)\cdot y)=x\cdot f_*(y)
 \qquad (x\in K^{\bullet}(X)_{\mathbb Q},\ y\in K^{\bullet}(X')_{\mathbb Q}).
\end{equation}


%%%%%%%%%%  scan idx 553  |  folio 540  %%%%%%%%%%
\SGAPrintPage{540}%

Passing to the associated graded groups in (3.2.1) and (3.2.2), one obtains a homomorphism of degree \(-r\):
\sgaStableTarget{sga6.e10.eq.3-2-3}\begin{equation}\tag{3.2.3}
 f_*:\Gro(X')_{\mathbb Q}\longrightarrow\Gro(X)_{\mathbb Q}
\end{equation}
and, by (1.3.2), a graded projection formula
\[
 f_*\bigl(f^*(x)\cdot y\bigr)=x\cdot f_*(y)
 \qquad (x\in\Gro(X)_{\mathbb Q},\ y\in\Gro(X')_{\mathbb Q}).
\]

\sgaStableTarget{sga6.e10.sec.3.3}\subsection*{3.3.}
Finally, consider a morphism \(f:X'\to X\) of finite Tor-dimension between noetherian schemes. Then it is known (Exp. IV) that \(f\) induces a homomorphism
\[
 f^*:K_\bullet(X)\longrightarrow K_\bullet(X')
\]
such that
\[
 f^*(\cl_\bullet(\underline{F}))=
 \sum_i(-1)^i\cl_\bullet\bigl(\underline{\Tor}_i^{\underline{O}_X}(\underline{F},\underline{O}_{X'})\bigr)
\]
for every coherent \(\underline{O}_X\)-module \(\underline{F}\).

The \(\underline{O}_X\)-modules \(\underline{\Tor}_i^{\underline{O}_X}(\underline{F},\underline{O}_{X'})\) are concentrated on \(f^{-1}(\Supp(\underline{F}))\), and when the fibers of \(f\) are everywhere of dimension \(\leq d\), it is known (EGA IV 5.6.7) that
\[
 \dim(f^{-1}(\Supp(\underline{F})))\leq \dim(\Supp(\underline{F}))+d.
\]
It follows that \(f\) induces, for every \(k\in\mathbb Z\), a homomorphism
\[
 f^*:\Filt_k(X)\longrightarrow\Filt_{k+d}(X'),
\]
whence a homomorphism of degree \(d\) of graded abelian groups


%%%%%%%%%%  scan idx 554  |  folio 541  %%%%%%%%%%
\SGAPrintPage{541}%

\sgaStableTarget{sga6.e10.eq.3-3-1}\begin{equation}\tag{3.3.1}
 f^*:\Grb(X)\longrightarrow\Grb(X').
\end{equation}

\sgaStableTarget{sga6.e10.sec.3.4}\subsection*{3.4.}
Finally, suppose that the morphisms (3.2.3) and (3.3.1) are both defined for a morphism \(f:X'\to X\) such that \(r=d\). This is the case, for example, when \(X\) is quasi-compact with an ample \(\underline{O}_X\)-module and \(X'=P_X(\underline{E})\), where \(\underline{E}\) is a locally free \(\underline{O}_X\)-module of rank \(r+1\). Then, for every \(x\in\Filt^i(X')_{\mathbb Q}\) and \(y\in\Filt_j(X)_{\mathbb Q}\), the projection formula in the derived category gives the projection formula
\[
 f_*\bigl(x\cdot f^*(y)\bigr)=f_*(x)\cdot y\in\Filt_{j+r}(X)_{\mathbb Q}.
\]
When \(X\) is noetherian, 1.3.2 yields the corresponding graded formula.

\sgaStableTarget{sga6.e10.sec.4}\section*{4. Intersection numbers}

\sgaStableTarget{sga6.e10.sec.4.1}\subsection*{4.1.}
Consider a proper scheme \(X\) over a field \(k\). Composing the product (3.1.1)
\[
 \Gr^i(X)\times\Gr_i(X)\longrightarrow \Gr_0(X)=\Filt_0(X)
 \qquad (i\in\mathbb Z)
\]
with the Euler--Poincaré characteristic
\[
 \chi:\Filt_0(X)\longrightarrow\mathbb Z,
\]
gives a pairing called the ``\emph{intersection number}'':
\sgaStableTarget{sga6.e10.eq.4-1-1}\begin{equation}\tag{4.1.1}
 \Gr^i(X)\times\Gr_i(X)\longrightarrow \mathbb Z,
\end{equation}


%%%%%%%%%%  scan idx 555  |  folio 542  %%%%%%%%%%
\SGAPrintPage{542}%

which we shall denote by
\[
 (x,y)\rightsquigarrow \langle x,y\rangle\in\mathbb Z
 \qquad (i\in\mathbb Z,\ x\in\Gr^i(X),\ y\in\Gr_i(X)).
\]

\sgaStableTarget{sga6.e10.ex.4-2}\textbf{Example 4.2.} Suppose that the scheme \(X\) has dimension \(n\). Take an element \(y\in\Gr_n(X)\), for example \(y=\cl_\bullet(\underline{O}_Z)\), where \(Z\) is an integral subscheme of dimension \(n\) in \(X\). Also take a homogeneous polynomial of weight \(n\)
\[
 P\in\mathbb Z[T_1,\ldots,T_n],
\]
where \(T_i\) is assigned weight \(i\) \((i=1,\ldots,n)\), for example a monomial \(T_{i_1}\cdot T_{i_2}\cdots T_{i_k}\) with \(i_1+i_2+\cdots+i_k=n\). Substituting in \(P\), for each variable \(T_i\) \((i=1,\ldots,n)\), the Chern class
\[
 c_i(x)=\cl^{\bullet}(\gamma^i(x-\epsilon(x)))\in\Gr^i(X)
 \qquad(x\in K^{\bullet}(X)),
\]
gives an element
\[
 P(c_1(x),\ldots,c_n(x))\in\Gr^n(X).
\]
The integer
\[
 P_y(x)=\bigl\langle P(c_1(x),\ldots,c_n(x)),y\bigr\rangle.
\]
is called the \emph{Chern number of} \(x\) \emph{relative to} \(y\).
When \(y=\cl_\bullet(\underline{O}_X)\), the reference to \(y\) is omitted from this terminology.

\sgaStableTarget{sga6.e10.sec.4.3}\subsection*{4.3.}
Without any hypothesis on the dimension of \(X\), consider a sequence \(\underline{L}_1,\ldots,\underline{L}_d\) of invertible \(\underline{O}_X\)-modules and a coherent \(\underline{O}_X\)-module \(\underline{F}\) such that \(\dim(\Supp(\underline{F}))\leq d\). Put \(L_i=\cl^{\bullet}(\underline L_i)\) \((i=1,\ldots,d)\) and \(F=\cl_\bullet(\underline F)\).


%%%%%%%%%%  scan idx 556  |  folio 543  %%%%%%%%%%
\SGAPrintPage{543}%

\sgaStableTarget{sga6.e10.prop.4-3-1}\textbf{Proposition 4.3.1.} \uline{With the above notation, for every sequence} \(\alpha_1,\ldots,\alpha_d\) \uline{of nonnegative integers with} \(\sum_{i=1}^{d}\alpha_i=d\)\uline{, the coefficient of}
\[
 \binom{n_1}{\alpha_1}\cdots\binom{n_d}{\alpha_d}
\]
\uline{in the Snapper polynomial} \(\chi(L_1^{n_1}\cdots L_d^{n_d}F)\) \uline{is equal to}
\sgaStableTarget{sga6.e10.eq.4-3-2}\begin{equation}\tag{4.3.2}
 \bigl\langle c_1(L_1)^{\alpha_1}\cdots c_1(L_d)^{\alpha_d},\lvert\underline{F}\rvert\bigr\rangle,
\end{equation}
\uline{where} \(\lvert\underline{F}\rvert\) \uline{is the class of} \(\underline{F}\) \uline{in} \(\Gr_d(X)\)\uline{. Thus, in the ordinary expansion of this polynomial in monomials in the} \(n_i\) \((i=1,\ldots,d)\)\uline{, the coefficient of the monomial} \(n_1^{\alpha_1}\cdots n_d^{\alpha_d}\) \uline{is equal to}
\[
 (1/(\alpha_1! \cdots \alpha_d!))
 \bigl\langle c_1(L_1)^{\alpha_1}\cdots c_1(L_d)^{\alpha_d},\lvert\underline{F}\rvert\bigr\rangle.
\]

In particular, one immediately obtains the following

\sgaStableTarget{sga6.e10.cor.4-3-3}\textbf{Corollary 4.3.3.} \uline{The coefficient of the monomial} \(n_1n_2\cdots n_d\) \uline{in the Snapper polynomial} \(\chi(L_1^{n_1}\cdots L_d^{n_d}F)\) \uline{is equal to}
\sgaStableTarget{sga6.e10.eq.4-3-4}\begin{equation}\tag{4.3.4}
 \bigl\langle c_1(L_1)\cdots c_1(L_d),\lvert\underline{F}\rvert\bigr\rangle
 =
 (c_1(L_1)\cdots c_1(L_d),\lvert\underline{F}\rvert).
\end{equation}

\emph{Proof}: Put \(L_i=1+M_i\) for every \(i=1,\ldots,d\). The definitions immediately show that the integer (4.3.2) is equal to \(\chi(M_1^{\alpha_1}\cdots M_d^{\alpha_d}F)\). Substituting \(A=K^{\bullet}(X)/\Filt^{d+1}(X)\), \(I=\Filt^1(X)/\Filt^{d+1}(X)\) in (2.1.4), formula (2.1.5) gives (taking account of the fact that \(M_i\in\Filt^1(X)\) for every \(i=1,\ldots,d\)) the congruence
\sgaStableTarget{sga6.e10.eq.4-3-5}\begin{equation}\tag{4.3.5}
 L_1^{n_1}\cdots L_d^{n_d}
 \equiv
 \sum_{\substack{\alpha_1,\ldots,\alpha_d\geq0\\ \alpha_1+\cdots+\alpha_d\leq d}}
 \binom{n_1}{\alpha_1}\cdots\binom{n_d}{\alpha_d}
 M_1^{\alpha_1}\cdots M_d^{\alpha_d}
 \quad\mathrm{mod.}\Filt^{d+1}(X).
\end{equation}
Thus the coefficient of \(\binom{n_1}{\alpha_1}\cdots\binom{n_d}{\alpha_d}\) in the polynomial \(\chi(L_1^{n_1}\cdots L_d^{n_d}F)\) is equal to \(\chi(M_1^{\alpha_1}\cdots M_d^{\alpha_d}F)\), Q.E.D.

